\section{Larger collections}
\subsection{Total cost}
To account for one complete request, let $T_{\rm shared}$ include query embedding,
normalization, routing, and any common filtering. If $T_a^{\rm local}$ is the local cost for
A, B, or C derived earlier, including its score table, we have
\begin{equation}
T_a^{\rm request}=T_{\rm shared}+T_a^{\rm local}+T_{\rm final},
\label{eq:total}
\end{equation}
where $T_{\rm final}$ contains output work and any enabled reranking. All of these stages
belong to the same request, so the equation does not justify adding unrelated stage medians
or adding the real encoder's measured time to constructed vectors it never produced.

For memory, we count the stored index separately from the temporary memory used during
a query. Write $M_0$ for memory allocated to document rows, routing data, and common containers.
Let $D_C$ be C's directory allocation, $Q_B$ B's extra masks and queue, and $Q_C$ C's extra
key queue. All methods also use the common temporary query memory $Q_0$.
\begin{center}
\begin{tabular}{@{}lll@{}}\toprule
Method & Stored index memory & Temporary query memory\\\midrule
A & $M_0$ & $Q_0$\\
B & $M_0+M_{\rm planes,allocated}$ & $Q_0+Q_B$\\
C & $M_0+D_C$ & $Q_0+Q_C$\\\bottomrule
\end{tabular}
\end{center}
The common data structures can reserve different amounts of unused space, so their actual
allocations can differ. A process RAM limit must also count memory used to run the program
and other data still in memory. We count each allocation once, even if both process and GPU
reports include it. A limit during index building must also include temporary input and output copies.

\subsection{Memory requirements}
At $N=10^9$, $d=256$, and eight-byte IDs, packed codes occupy 32 GB and IDs occupy 8 GB,
so these two arrays alone exceed a 32 GB budget. B also keeps a second sign representation,
adding at least 32 GB before per-cluster padding. Changing how many clusters we open
does not reduce these stored bytes.

\begin{center}
\begin{tabular}{@{}lrrl@{}}\toprule
Layout at 1B rows & Minimum data & Planned query RAM & 64 GB budget\\\midrule
A, direct prefix route & 40 GB & about 49 GB & Predicted to fit\\
C, global 13-bit key & 40 GB & about 49 GB & Predicted to fit\\
B, global cluster & 72 GB & about 83 GB & Rejected\\
B, prefix route, large leaf & 72 GB & about 81 GB & Rejected\\\bottomrule
\end{tabular}
\end{center}
These estimates allow for reserved space and data structures, plus a configurable 1 GB
reserve for running the program. That reserve is an assumption, not a memory measurement. In their current in-memory layouts,
all three fail 32 GB. All listed planning cases fit 1 TB, but a machine with 1 TB need not
have this machine's memory performance. If we compressed the IDs or used another storage
architecture, we would need to calculate memory again; relabeling the same index would not
change the result.

\subsection{Fixed bit positions}
If we keep thirteen strong coordinates fixed under the earlier independent strong-coordinate
model, the expected ideal group contains $10^9/2^{13}\approx122{,}070$ rows, far more than
the requested 100. The earlier probability calculation accounts for the small chance that fewer than 100
rows match, under these assumptions. It makes no claim about ordinary text embeddings.

A can route directly to this group and scan it, while C can look up its key and score the
same rows. Under these assumptions, the model fitted to measured times predicts roughly 3 ms for both. B can also take this
route with a large leaf, giving a similar time with more stored data. We should not force B
through a global billion-row bitmap just to make C appear better. In this case, A and C are closely
related ways to obtain the same candidate set, and the model cannot reliably distinguish a gap of a few percent.

\subsection{Changing bit positions}
If the important coordinates change with each query, a fixed prefix may contain few
of the coordinates that determine the ranking, while B can choose its split coordinates from the query. We preserve the
constructed path by keeping thirteen useful cuts and scaling the leaf size to 160,000.
At the last split, the masks occupy approximately
\[
(13+2)\ceil{10^9/64}\times8=1.875\ \text{GB}.
\]
The factor 15 counts masks on this path and queued alternatives held at the same time,
including the parent while its children are created. This count applies to the specified
path; other sequences of splits can need different amounts of temporary memory.

The model for the complete C++ search call predicts about 0.10 s for this global B path and about 25 s
for a global scan. If we double B's variable cost, we get about 0.20 s, which shows how the prediction changes
when that assumed cost changes; it is not a confidence interval. No full billion-document query was measured.

We also need to give A the benefit of routing in this comparison. If it scores fraction $f$
of the collection, its approximate cost is $T_{\rm route}+fNc_{\rm scan}$. Even with routing
time ignored, A must score less than about 0.41\% of the collection to beat the 0.10 s B
prediction, while still meeting the recall target at that fraction. Scanning can win if
routing cheaply produces a small, accurate candidate set; if it cannot, the query-dependent
B path is worth testing.

\fig{billion-comparison}{A conditional comparison with sufficient RAM. B is one prescribed operating point. A gets free routing in this line but must still attain the required recall at its scanned fraction. These are forecasts, not measured billion-document timings.}{.9}

For C, the short key has to capture important coordinates. In the constructed changing case,
a fixed 13-coordinate key misses all thirteen strong coordinates with probability about
49.94\%. A lookup that finds documents can therefore still leave almost all the scoring work to do.
The number of documents found by a key alone cannot establish faster search or better recall.

\subsection{Conclusions}
For the real embeddings measured here, the scan remains the supported practical baseline.
The constructed comparisons show how that result depends on the workload: when the strong coordinates
stay in the fixed prefix, direct routing and key lookup both avoid most scores; with changing
strong coordinates, B's query-dependent splits produced a large measured advantage. Each
conclusion applies to its stated workload and tested parameter space.

We can use the formulas to explain these outcomes and suggest settings, but they do not
prove that every possible IVF layout loses, that Matryoshka creates the constructed
distribution, or that a fitted cost per operation remains exact at a billion rows. To test
the favorable case on a real representation, the next useful experiment would check whether a few coordinates determine the ranking
and whether routing finds the right documents, under the same recall and RAM limits.
